\documentclass[11pt,a4paper]{article}
\usepackage[utf8]{inputenc}
\usepackage[T1]{fontenc}
\usepackage[french]{babel}
\usepackage{geometry}
\geometry{margin=2.5cm}
\usepackage{amsmath,amssymb}
\usepackage{graphicx}
\usepackage{listings}
\usepackage{xcolor}
\usepackage{hyperref}
\hypersetup{colorlinks=true, linkcolor=blue!60!black, urlcolor=blue!60!black}

\lstset{
  language=Python,
  basicstyle=\ttfamily\small,
  keywordstyle=\color{blue!70!black},
  commentstyle=\color{green!50!black},
  stringstyle=\color{red!60!black},
  frame=single,
  breaklines=true,
  captionpos=b,
}

\title{Sliceur non planaire pour impression FDM\\sur bras robotique à 6 degrés de liberté}
\author{}
\date{\today}

\begin{document}
\maketitle
\tableofcontents

% ======================================================================
\section{Introduction}
% ======================================================================

L'impression 3D par dépôt de matière fondue (FDM) est presque
exclusivement réalisée sur des machines cartésiennes à 3 axes: la
buse reste verticale et chaque couche est un plan horizontal parallèle
au plateau. Cette contrainte découle du logiciel, pas de la physique:
rien n'empêche mécaniquement de déposer du filament avec une buse
inclinée, sur une surface courbe, tant que la cinématique le permet.
Un bras robotique à 6 degrés de liberté (DDL) possède précisément les
six degrés nécessaires pour contrôler, en chaque point du trajet, la
position \emph{et} l'orientation de l'outil.

Ce projet développe un \emph{sliceur non planaire} en Python: un
logiciel qui transforme un modèle mécanique (pièce composite à fibres
continues, moule support inférieur) en trajectoires d'impression 3D
complètes --- position, orientation et codes de déplacement --- prêtes
à être exécutées en simulation sur RoboDK. Le code s'appuie sur
l'écosystème scientifique Python (NumPy, SciPy) et sur la bibliothèque
de visualisation et de traitement de maillages PyVista (VTK).

L'architecture logicielle repose sur trois classes:
\begin{itemize}
  \item \texttt{SlicedObject}: classe parente, une couche de surface
        quelconque et les opérations génériques (champ de distance au
        bord, surfaces décalées, contours, construction et export de
        trajectoire);
  \item \texttt{SlicedPiece}: couche de la pièce composite à imprimer,
        avec prise en compte des chemins de fibres et de leur direction
        principale;
  \item \texttt{SlicedMold}: couche du moule de support généré
        automatiquement sous les surplombs de la pièce, avec sa
        logique propre de volume d'ombre, d'infill et de connexions.
\end{itemize}
Le présent rapport explique les choix méthodologiques et
d'implémentation du projet selon la structure suivante: les slicers non
planaires existants (section~\ref{sec:related}), le dispositif de
fabrication (section~\ref{sec:setup}), les exigences imposées au sliceur
(section~\ref{sec:requirements}), la structure de calcul
(section~\ref{sec:structure}), l'implémentation (section~\ref{sec:impl})
et une discussion des résultats et limites (section~\ref{sec:discussion}).

% ======================================================================
\section{Slicers non planaires existants}\label{sec:related}
% ======================================================================

Les slicers conventionnels (Cura, PrusaSlicer, Slic3r) découpent le
modèle STL en \emph{couches planes} équidistantes: intersection du
maillage avec une famille de plans $z = z_0 + k\,h$, remplissage 2D
(infill à $\pm45^\circ$, périmètres), génération de G-code à 3 axes.
Ce paradigme est robuste mais impose deux limitations majeures: la
staircase effect (escaliers) sur les surfaces inclinées, et
l'impossibilité de déposer sur des surfaces non développables.

Plusieurs approches ont été proposées pour lever ces limites:
\begin{itemize}
  \item \textbf{Découpe non plane par surfaces de niveau}~: au lieu de
        plans, la pièce est découpée par des surfaces courbes épousant
        la géométrie (par exemple des isosurfaces d'un champ scalaire).
        C'est la voie retenue ici: les couches sont des isosurfaces
        $q = c$ d'un champ de \og hauteur de couche \fg{} défini sur le
        maillage volumique (section~\ref{sec:structure}).
  \item \textbf{Impression conformelle sur surfaces existantes}
        (dépôt sur un substrat courbe, hybride impression/usinage):
        la trajectoire est paramétrée par la surface du substrat, la
        buse restant normale à celui-ci.
  \item \textbf{Impression multi-axes}~: des travaux académiques
        (p.ex.\ curved layer slicing pour bras 6-DDL) montrent que
        l'orientation de la buse doit être continûment contrôlée le
        long du trajet, ce qui suppose un sliceur produisant une
        \emph{pose} complète $(R, t)$ et non plus des positions seules.
  \item \textbf{Supports conformes}~: les slicers classiques génèrent
        des supports en couches planes; la génération d'un moule
        épousant la surface inférieure de la pièce (volume d'ombre)
        est une extension non plane de la notion de support.
\end{itemize}

Le présent projet se distingue par trois points: (i) le découpage en
couches est \emph{non planaire par construction} (isosurfaces d'un
champ scalaire, calculées via VTK), (ii) l'infill de la pièce est
\emph{aligné sur les fibres} déposées dans la couche (renfort continu),
et (iii) chaque point de trajectoire est exporté avec son repère
outil complet, directement consommable par RoboDK pour un bras 6-DDL.

% ======================================================================
\section{Dispositif de fabrication}\label{sec:setup}
% ======================================================================

Le dispositif cible est une cellule d'impression robotisée:
\begin{itemize}
  \item \textbf{Bras robotique à 6 DDL} (par ex.\ bras industriel
        articulé), capable de positionner l'outil en position et en
        orientation arbitraires dans son espace de travail.
  \item \textbf{Tête de dépôt} montée en effecteur: extrudeuse FDM
        (buse, filament thermoplastique ou composite à fibres
        continues). La buse n'est \emph{pas} rotationnellement
        symétrique autour de son axe d'approche: son orientation propre
        (rotation autour de l'axe) a donc un sens mécanique.
  \item \textbf{RoboDK} pour la simulation hors ligne: le sliceur
        exporte des fichiers texte de poses $(x,y,z, R_{3\times3})$
        avec codes de segment, que RoboDK convertit en programme
        robot.
\end{itemize}

Le flux global est le suivant:
\begin{enumerate}
  \item Un modèle mécanique (maillage volumique + moule, issu d'une
        simulation par éléments finis --- fichier APDL parse par la
        bibliothèque \texttt{oamc}) est chargé.
  \item Le modèle est orienté: une matrice de rotation aligne la
        normale moyenne des couches sur l'axe $Z$ du monde
        (\texttt{SlicedPiece.get\_rotation\_matrix}, puis
        \texttt{SlicedMold.rotate\_model} et \texttt{lift\_model}).
  \item Le sliceur génère les couches de la pièce
        (\texttt{SlicedPiece}) et les couches du moule support
        (\texttt{SlicedMold}, à partir du volume d'ombre).
  \item Les trajectoires (segments + repères outil + codes) sont
        exportées vers des fichiers texte et importées dans RoboDK
        pour simulation et génération du programme robot.
\end{enumerate}

Paramètres de fabrication typiques utilisés dans le code: hauteur de
couche $h = 0{,}35$\,mm, diamètre de buse $d = 0{,}70$\,mm, distance
d'approche 1\,mm, pourcentage d'infill réglable (espacement
$d \times p^{-1}$, cf.\ \texttt{infill\_spacing}).

% ======================================================================
\section{Exigences pour le sliceur non planaire}\label{sec:requirements}
% ======================================================================

L'analyse du dispositif et des objectifs fonctionnels conduit au jeu
d'exigences suivant.

\subsection{Exigences géométriques}
\begin{itemize}
  \item \textbf{E1 --- couches non planes}~: les surfaces de couche
        suivent la géométrie de la pièce (isosurfaces d'un champ
        scalaire), pas des plans horizontaux.
  \item \textbf{E2 --- infill aligné sur les fibres}~: dans la pièce
        composite, les lignes d'infill doivent être orientées par
        rapport à la direction moyenne des fibres déposées (angle de
        $\pm45^\circ$ alterné par couche), les fibres étant chargées
        depuis des fichiers CSV.
  \item \textbf{E3 --- contour shells}~: chaque couche doit comporter
        des périmètres (shells) décalés de la frontière d'un multiple
        du diamètre de buse, sur la pièce comme sur le moule.
  \item \textbf{E4 --- moule automatique}~: le support sous les
        surplombs doit être généré automatiquement comme volume d'ombre
        de la pièce, puis découpé en couches, avec une infill en grille
        $XY$ et gestion des croisements.
\end{itemize}

\subsection{Exigences robotiques}
\begin{itemize}
  \item \textbf{E5 --- pose complète}~: chaque point de trajectoire doit
        être livré avec son repère orthonormé direct: axe $z$ = axe
        d'approche (normal à la surface, pointé vers la surface), axe
        $x$ = direction de parcours, axe $y = z \times x$. L'outil n'est
        pas symétrique en rotation: l'orientation doit être continue et
        déterministe pour éviter les collisions et les retournements.
  \item \textbf{E6 --- mouvements d'approche/retrait}~: chaque
        sous-trajet commence par une approche le long de la normale
        (non imprimée) et se termine par un retrait symétrique; les
        connexions trop longues ou interdites deviennent
        retrait + approche.
  \item \textbf{E7 --- codes de segment}~: les mouvements doivent être
        typés (approche, courbe, connexion, retrait) avec des codes
        distincts pour les passes \og pure \fg{} (shells) et
        \og composite \fg{} (infill), exploitables par le robot.
  \item \textbf{E8 --- export texte multi-couches}~: un fichier unique
        accumule toutes les couches, chaque ligne portant
        éventuellement l'identifiant de couche.
\end{itemize}

\subsection{Exigences logicielles et de performance}
\begin{itemize}
  \item \textbf{E9 --- architecture objet}~: factorisation du
        comportement commun (classe \texttt{SlicedObject}) avec
        spécialisation par type d'objet imprimé.
  \item \textbf{E10 --- performance}~: les opérations coûteuses
        (normales, arbres KD, localisateurs, requêtes plus proche
        cellule) doivent être calculées une fois, mises en cache et
        appelées de façon vectorisée (par lots), pas point par point.
  \item \textbf{E11 --- simplification contrôlée}~: les polylignes
        exportées doivent pouvoir être allégées (Ramer--Douglas--Peucker)
        sans jamais déplacer les extrémités ni fusionner des segments.
\end{itemize}

% ======================================================================
\section{Structure de calcul}\label{sec:structure}
% ======================================================================

Le pipeline de calcul se décompose en cinq étapes, toutes construites
sur PyVista/VTK.

\subsection{Découpage en couches (pièce)}
Sur le maillage volumique, un champ scalaire $q$ est défini tel que
chaque niveau entier de $q$ corresponde à une couche d'impression. La
fonction VTK de contour (surface de niveau) fournit alors la surface
de chaque couche:
\begin{equation}
  \mathcal{S}_k \;=\; \{\, \mathbf{p} \in \mathcal{V} : q(\mathbf{p}) = q_k \,\},
\end{equation}
où $\mathcal{V}$ est le maillage volumique. C'est la généralisation
directe du slicing plan (où $q = (z - z_0)/h$) à des couches courbes.

\subsection{Champ de distance au bord et contours}
Pour chaque surface de couche, on extrait les arêtes de frontière
(\texttt{extract\_feature\_edges}) et on construit un arbre KD
(\texttt{scipy.spatial.cKDTree}) sur les points de bord. Le champ
\texttt{dist\_to\_boundary} est la distance de chaque point de la
surface au bord le plus proche:
\begin{equation}
  d(\mathbf{p}) = \min_{\mathbf{b} \in \partial\mathcal{S}} \|\mathbf{p}-\mathbf{b}\|.
\end{equation}
Les contour shells sont les isolignes $d = c_i$ (par ex.\
$c_i \in \{0{,}5,\,1{,}0,\,\dots\}$\,mm), et la zone d'infill est la
sous-surface $d \ge c_{\max}$ obtenue par découpe scalaire
(\texttt{clip\_scalar}, avec re-triangulation le long de l'isovaleur).

\subsection{Infill}
\textbf{Pièce}: la direction moyenne des fibres est la première
composante principale des directions de segments (eigenvecteur
dominant de la matrice de second moment non centrée
$\mathbf{D}^\top\mathbf{D}$, cf.\ \S\ref{sec:impl-piece}). Les plans de
coupe sont normaux à une direction à $\pm45^\circ$ de cette direction
moyenne dans le plan tangent, alternés par couche; on découpe la zone
interne par tous les plans en un seul appel au filtre de contour
(niveaux multiples d'un champ scalaire $\mathbf{p}\cdot\mathbf{m}$),
puis on chaîne les segments en polylignes ordonnées.

\textbf{Moule}: grille $XY$ à espacement
$s = d_{\text{buse}} / p_{\text{infill}}$; les positions des plans
d'infill sont figées \emph{dès la couche zéro} et réutilisées pour
toutes les couches, afin que l'infill s'empile verticalement. À chaque
couche, la direction primaire ($X$ ou $Y$ selon parité) reste continue
et la direction secondaire est coupée aux intersections avec la
primaire (sinon le croisement serait imprimé deux fois).

\subsection{Assemblage de la trajectoire}
Les courbes (shells puis infill) sont recousues en une passe
connectée par un algorithme glouton de plus proche extrémité: à chaque
itération, on choisit l'extrémité libre la plus proche, on connecte,
et si la connexion dépasse une longueur maximale --- ou si elle est
interdite par un \emph{blocker} (croisement d'infill) --- on la
remplace par retrait + approche le long de la normale locale.

\subsection{Orientation et export}
Pour chaque point de trajectoire: normale locale $\mathbf{n}$ (cellule
la plus proche, de la surface pour la pièce, de la \emph{surface de la
pièce} pour le moule), tangente de parcours $\mathbf{t}$ (différence
centrale), puis repère orthonormé direct par Gram--Schmidt:
\begin{equation}
  \mathbf{z} = -\mathbf{n} \text{ (axe d'approche vers la surface)}, \quad
  \mathbf{x} = \frac{\mathbf{t} - (\mathbf{t}\cdot\mathbf{z})\,\mathbf{z}}{\|\cdot\|}, \quad
  \mathbf{y} = \mathbf{z}\times\mathbf{x}.
\end{equation}
Le tout est exporté en texte: $x\,y\,z$, les neuf composantes du
repère (colonnes de $R$), le type de segment et la couche.

% ======================================================================
\section{Implémentation}\label{sec:impl}
% ======================================================================

Le paquet \texttt{slicer} (disposition \texttt{src/}) compte trois
modules, ~2750 lignes au total.

\subsection{\texttt{SlicedObject}: classe parente (\texttt{SlicedObject.py})}
\label{sec:impl-object}

Responsable d'une couche générique, elle fournit:
\begin{itemize}
  \item \textbf{gestion de surface avec invalidation de caches}: la
        propriété \texttt{surface} réinitialise les caches (normales de
        cellules, champ de distance) à chaque réassignation --- nécessaire
        car \texttt{merge\_points} change l'indexation des points.
  \item \textbf{champ de distance au bord}
        (\texttt{build\_distance\_tree}): \texttt{cKDTree} sur les points
        d'arêtes, requête parallèle (\texttt{workers=-1}).
  \item \textbf{surface décalée} (\texttt{build\_offset\_surface}):
        \texttt{clip\_scalar} sur \texttt{dist\_to\_boundary}, la
        convention \texttt{invert} étant documentée explicitement dans
        le code (source classique d'inversion de résultat).
  \item \textbf{polylignes}: découpage de la connectivité \texttt{lines}
        de VTK en listes de tableaux NumPy
        (\texttt{\_polylines\_from\_polydata}), ordonnancement de
        segments via \texttt{vtkStripper} (\texttt{order\_segments}).
  \item \textbf{construction du chemin}
        (\texttt{build\_path\_single\_pass}): table d'extrémités
        $(2N \times 3)$, masque d'extrémités actives, distance
        vectorisée au plus proche voisin à chaque itération, normales
        d'extrémités calculées en un lot. Un \texttt{connection\_blocker}
        optionnel permet d'interdire des connexions (utilisé par le
        moule pour éviter de réimprimer un croisement).
        \texttt{build\_full\_printing\_path} enchaîne deux passes
        (pure/composite) avec des codes G-code-like distincts
        (approche 4/5, fin de connexion 1/3, continuation 0/2).
  \item \textbf{orientation et export}:
        \texttt{compute\_path\_normals} (requête par lots de la cellule
        la plus proche), \texttt{compute\_path\_tangents} (différences
        centrales, unilatères aux extrémités),
        \texttt{build\_orthonormal\_frames} (Gram--Schmidt vectorisé,
        repli arbitraire quand tangente $\parallel$ normale ---
        approche/retrait),
        \texttt{simplify\_polyline} (RDP itératif, en 3D),
        \texttt{export\_path\_to\_txt} / \texttt{export\_printing\_path}
        (écriture par segment via \texttt{np.savetxt}, mode
        \texttt{w}/\texttt{a} pour l'accumulation multi-couches,
        \texttt{drop\_shared\_points} pour éviter les points doublons
        aux jonctions).
\end{itemize}

\subsection{\texttt{SlicedPiece}: la pièce (\texttt{SlicedPiece.py})}
\label{sec:impl-piece}

Spécialise la couche pour la pièce composite:
\begin{itemize}
  \item \textbf{fibres}: chargement CSV par couche
        (\texttt{load\_fiber\_paths}); direction moyenne = vecteur
        propre dominant de $\mathbf{D}^\top\mathbf{D}$ (PCA non centrée
        --- un centrage donnerait l'axe de dispersion, perpendiculaire
        au faisceau!), signe rendu déterministe.
  \item \textbf{plans de coupe}
        (\texttt{generate\_slicing\_planes}): direction à
        $\pm45^\circ$ de la direction des fibres dans le plan tangent,
        normales de plans $\mathbf{m} = \mathbf{n} \times \mathbf{d}$,
        positions réparties sur l'étendue projetée avec marge d'un
        demi-pas.
  \item \textbf{coupe en une passe} (\texttt{\_cut\_with\_planes}): les
        plans étant parallèles, ils sont des niveaux du champ
        $\mathbf{p}\cdot\mathbf{m}$; un seul appel au filtre de contour
        remplace $n$ appels à \texttt{slice}. Le chaînage
        (\texttt{\_chain\_segments}) reconstruit les polylignes en
        parcourant le graphe de segments (degrés $\le 2$), avec repli
        sur une méthode plus lente si le graphe contient des jonctions.
  \item \textbf{préparation du modèle}: extraction des surfaces de
        niveau (\texttt{get\_printing\_layers},
        \texttt{slice\_model}, \texttt{from\_model}), rotation moyenne
        des normales vers $+Z$ (\texttt{get\_rotation\_matrix} via
        \texttt{scipy Rotation}).
  \item \textbf{chemin complet} (\texttt{build\_printing\_path}):
        fusion des points coïncidents d'abord (sinon la détection
        d'arêtes de frontière est faussée), puis zone interne
        $\Rightarrow$ infill aligné fibres, isolignes de distance
        $\Rightarrow$ shells, recousus par
        \texttt{build\_full\_printing\_path}.
\end{itemize}

\subsection{\texttt{SlicedMold}: le moule (\texttt{SlicedMold.py})}
\label{sec:impl-mold}

Le module le plus riche (~1590 lignes). Points clés:
\begin{itemize}
  \item \textbf{normales d'approche redéfinies}
        (\texttt{get\_local\_normals}): les approches/retraits du moule
        sont perpendiculaires à la \emph{surface de la pièce} (contre
        laquelle le support est imprimé), pas à la couche du moule
        elle-même.
  \item \textbf{volumes d'ombre}: trois constructions ---
        projection verticale de la surface inférieure
        (\texttt{get\_shadow\_volume}), cône de rayons lancés depuis
        les faces vers le bas (\texttt{get\_shadow\_volume\_cone}, avec
        pré-filtres vectorisés et localisateur unique
        \texttt{\_RayCaster} autour de \texttt{vtkStaticCellLocator}),
        et projection verticale simple
        (\texttt{get\_shadow\_volume\_vertical}). Le contour d'ombre est
        le convexe hull 2D de la projection (\texttt{compute\_shadow\_contour}).
  \item \textbf{découpage} (\texttt{slice\_shadow\_volume}): champ
        $q = (Z - z_0)/h$ sur le volume; pour chaque niveau entier:
        isosurface, triangulation, champ \texttt{dist\_to\_part}
        (minimum de la distance implicite à la pièce triangulée et de
        la hauteur au plan), shells isolignes de
        \texttt{dist\_part}, zone interne par \texttt{clip\_scalar},
        infill $X$ et $Y$ aux niveaux \emph{figés depuis la couche
        zéro} (pour un empilement vertical exact).
  \item \textbf{traitement de l'infill}~: intersections 2D
        vectorisées (\texttt{\_segment\_intersections\_2d}), découpe
        des lignes secondaires aux croisements
        (\texttt{cut\_infill\_at\_intersections\_fast}), longueur de
        connexion maximale dérivée des paramètres
        (\texttt{compute\_max\_connection\_length}), réduction d'une
        ligne d'infill à ses extrémités
        (\texttt{collapse\_infill\_line}).
  \item \textbf{chemin de couche} (\texttt{build\_support\_path}):
        tri zigzag des lignes, parité de couche $\Rightarrow$ direction
        primaire $X$ ou $Y$ continue, secondaire coupée; deux passes
        (shells puis infill) recousues par
        \texttt{build\_path\_single\_pass} avec un
        \texttt{connection\_blocker} interdisant les connexions qui
        traverseraient une ligne primaire.
  \item \textbf{export multi-couches}
        (\texttt{export\_layers\_to\_txt}).
\end{itemize}

\subsection{Exemples d'utilisation}
Le dépôt fournit deux scripts dans \texttt{examples/}:
\texttt{print\_layers.py} (pièce: parsing APDL, matériaux matrice/fibre,
rotation et élévation du modèle, chargement des fibres, découpage,
trajectoires) et \texttt{mold\_generation.py} (moule: mêmes étapes
amont puis génération du volume d'ombre et découpage). Les chemins de
fibres CSV d'exemple (\texttt{examples/fibers/layer\_*/}) alimentent
l'infill aligné.

% ======================================================================
\section{Discussion}\label{sec:discussion}
% ======================================================================

\subsection{Points forts}
\begin{itemize}
  \item \textbf{Généralité du découpage}: le passage par des
        isosurfaces rend le slicing non planaire aussi simple que le
        slicing plan; le même code sert pour la pièce et le moule.
  \item \textbf{Poses complètes et cohérentes}: le couple
        normale--tangente, orthogonalisé en repère direct, donne au
        robot une orientation continue et mécaniquement significative
        --- indispensable pour un outil non symétrique.
  \item \textbf{Qualité des trajectoires}: approches/retraits normaux à
        la surface, connexions bornées, croisements d'infill jamais
        imprimés deux fois, plans d'infill du moule empilés
        verticalement.
  \item \textbf{Performance}: l'ensemble des opérations critiques est
        mis en cache ou vectorisé (normales une seule fois, arbres KD
        uniques, découpes multi-niveaux en une passe, requêtes par
        lots), ce qui rend le traitement de maillages FEM complets
        praticable en Python.
\end{itemize}

\subsection{Limites et travaux futurs}
\begin{itemize}
  \item \textbf{Collision et cinématique}~: le sliceur ignore l'espace
        de travail du bras et les auto-collisions; RoboDK les détecte
        en simulation mais pas le sliceur en amont. Un module de
        vérification de faisabilité (joint limits, proximité
        pièce/bras) serait la suite naturelle.
  \item \textbf{Continuité d'orientation inter-couches}~: les tangentes
        sont calculées par segment; une légère discontinuité de
        rotation propre peut apparaître aux jonctions de segments
        (sans effet sur la position).
  \item \textbf{Chargement des fibres}~: une seule couche de fibres est
        chargée à la fois (TODO documenté dans le code); le
        couplage complet avec la génération charge-based de
        \texttt{oamc} reste à faire.
  \item \textbf{Validation expérimentale}~: le pipeline est validé en
        simulation RoboDK; l'impression réelle (dynamique d'extrusion
        sur surface courbe, adhésion sur moule) reste à caractériser.
\end{itemize}

\subsection{Conclusion}
Le projet démontre qu'un sliceur non planaire complet --- couches
courbes, infill aligné sur les fibres, moule support automatique,
poses 6-DDL exportables --- peut être construit de façon concise en
s'appuyant sur PyVista/VTK et NumPy, avec une architecture à trois
classes qui sépare proprement ce qui est générique (géométrie,
assemblage, orientation, export) de ce qui est spécifique à l'objet
imprimé (pièce composite ou moule). Les prochaines étapes naturelles
sont la vérification cinématique en amont, le couplage complet avec la
génération de fibres et la validation expérimentale sur le bras.

\end{document}
